\clearpage
\section{Der Weg zu den Beweistabellen}
\label{cuts:correspondence}

Die kurzen Beweise dieses Hefts erläutern die folgenden Aussagen aus
Band~19. Die verlinkten \href{_B19-cuts-proofs.pdf\#cuts.proofs}{Beweistabellen}
führen ihre einzelnen Ableitungsschritte aus. Aussage und Nummer stimmen
mit dem Fachband überein.

\begin{center}
\small\renewcommand{\arraystretch}{1.25}
\begin{tabularx}{\linewidth}{@{}Xl@{}}
\toprule
Schritt & Aussage in Band 19\\
\midrule
Schnittbedingungen & \href{_B19-cuts-proofs.pdf\#cuts.proof.DedekindCutPredicate}{Def.~19.2.1.1}\\
Schnittkriterium & \href{_B19-cuts-proofs.pdf\#cuts.proof.DedekindCutCharacterization}{Th.~19.2.1.1}\\
Rationale Hauptschnitte & \href{_B19-cuts-proofs.pdf\#cuts.proof.RealPrincipalCutIsCut}{Th.~19.2.2.1}\\
Vergleichbarkeit & \href{_B19-cuts-proofs.pdf\#cuts.proof.RealCutInclusionComparable}{Th.~19.3.1.2}\\
Ordnungstreue & \href{_B19-cuts-proofs.pdf\#cuts.proof.RationalToRealEmbeddingOrder}{Th.~19.3.1.4}\\
Injektive Einbettung & \href{_B19-cuts-proofs.pdf\#cuts.proof.RationalToRealEmbeddingInjective}{Th.~19.3.1.5}\\
Vereinigung ist ein Schnitt & \href{_B19-cuts-proofs.pdf\#cuts.proof.RealBoundedCutUnionIsCut}{Th.~19.3.3.1}\\
Kleinste obere Schranke & \href{_B19-cuts-proofs.pdf\#cuts.proof.RealLeastUpperBoundProperty}{Th.~19.3.3.2}\\
Supremum als Vereinigung & \href{_B19-cuts-proofs.pdf\#cuts.proof.RealCutSupremumUnion}{Th.~19.3.3.3}\\
\bottomrule
\end{tabularx}
\end{center}

\textbf{Derselbe Gedanke in zwei Darstellungen.}
Beim Vergleich rationaler Hauptschnitte verwendet der Lesebeweis den
Randwert selbst als Gegenzeugen. Die Tabellen wählen an dieser Stelle
zum Teil eine rationale Zwischenzahl. Beide Wege begründen denselben
Ordnungsvergleich.

Beim Supremum bleibt die Reihenfolge entscheidend: Zuerst ist die
Vereinigung als Schnitt nachzuweisen; erst danach darf sie als reelle
obere Schranke auftreten. Die beiden Hilfsableitungen zur oberen Schranke
und zur Kleinsteigenschaft enthalten diese Voraussetzung ausdrücklich.

Die Beispiele und Übungen zeigen konkrete Anwendungen. Ihre Rechnungen
stehen im jeweiligen Heft. Das \href{_B19-cuts-exercises.pdf\#cuts.exercises}{Übungsheft}
endet mit einem selbst geführten Beweis der Supremumsformel.
